\newpage
\markboth{Hinweise zur Nutzung von KI}{Hinweise zur Nutzung von KI}
\chapter*{Hinweise zur Nutzung von KI}

Die Verwendung sämtlicher im Rahmen dieser Abschlussarbeit eingesetzten KI-Tools ist an dieser Stelle transparent und nachvollziehbar zu dokumentieren. Dabei sind die verwendeten KI-Tools konkret zu benennen sowie Art, Zweck und Umfang der Nutzung zu beschreiben.

Zur Orientierung sind nachfolgend mögliche Einsatzbereiche von KI aufgeführt. Auf die jeweils zutreffenden Bereiche ist in dieser Erklärung einzugehen:
\begin{itemize}
  \item Ideenfindung (Brainstorming) und Strukturierung
  \item Recherche, Zusammenfassung oder Verwaltung von Literatur
  \item Datenaufbereitung und -analyse
  \item Programmierung und technische Umsetzung
  \begin{itemize}
    \item Vollständige Erstellung von Code
    \item Unterstützung bei der Optimierung von Code
  \end{itemize}
  \item Generierung oder sprachliche Überarbeitung von Texten sowie Übersetzung
  \item Visualisierung und Gestaltung
  \item Qualitätssicherung und Review
  \item Sonstiges
\end{itemize}

Anschließend ist die Verantwortung für die unter Verwendung von KI-Tools entstandenen Inhalte zu bestätigen. Hierfür kann beispielsweise folgende Formulierung verwendet werden:

\enquote{Nach der Verwendung der genannten KI-Tools hat der Autor/die Autorin die daraus entstandenen Inhalte eigenständig überprüft und, soweit erforderlich, bearbeitet und angepasst. Der Autor/die Autorin übernimmt die volle Verantwortung für den Inhalt und die Richtigkeit der vorliegenden Arbeit.}
